\documentclass[journal]{IEEEtran}

% ---- Packages ----
\usepackage[utf8]{inputenc}
\usepackage[T1]{fontenc}
\usepackage{amsmath,amssymb,amsfonts,amsthm}
\usepackage{algorithmic}
\usepackage{algorithm}
\usepackage{graphicx}
\usepackage{textcomp}
\usepackage{xcolor}
\usepackage{booktabs}
\usepackage{multirow}
\usepackage{hyperref}
\usepackage{url}
\usepackage{cite}

\hypersetup{
  colorlinks=true,
  linkcolor=black,
  citecolor=black,
  urlcolor=blue
}

% ---- Theorem environments ----
\newtheorem{theorem}{Theorem}
\newtheorem{lemma}[theorem]{Lemma}
\newtheorem{proposition}[theorem]{Proposition}
\newtheorem{corollary}[theorem]{Corollary}
\newtheorem{definition}[theorem]{Definition}
\newtheorem{remark}[theorem]{Remark}
\newtheorem{example}[theorem]{Example}
\newtheorem{axiom}{Axiom}

% ---- Custom commands ----
\newcommand{\R}{\mathbb{R}}
\newcommand{\E}{\mathbb{E}}
\newcommand{\N}{\mathbb{N}}
\newcommand{\norm}[1]{\left\lVert #1 \right\rVert}
\newcommand{\ip}[2]{\langle #1, #2 \rangle}
\newcommand{\KL}{\mathrm{KL}}
\newcommand{\Hcal}{\mathcal{H}}
\newcommand{\Mcal}{\mathcal{M}}
\newcommand{\Scal}{\mathcal{S}}
\newcommand{\Bcal}{\mathcal{B}}
\newcommand{\Pcal}{\mathcal{P}}
\DeclareMathOperator*{\argmin}{arg\,min}
\DeclareMathOperator*{\argmax}{arg\,max}
\DeclareMathOperator{\tr}{tr}
\DeclareMathOperator{\rank}{rank}
\DeclareMathOperator{\Vol}{Vol}
\DeclareMathOperator{\diam}{diam}

\begin{document}

% ---- Title ----
\title{Geometric Aesthetics: Beauty as Compressibility\\in Perceptual Eigenspaces}

\author{Andrew~H.~Bond
\thanks{Department of Computer Engineering, San Jose State University.
E-mail: andrew.bond@sjsu.edu}}

\markboth{Geometric Aesthetics}{Bond: Beauty as Compressibility in Perceptual Eigenspaces}

\maketitle

% =====================================================================
% ABSTRACT
% =====================================================================
\begin{abstract}
We develop a mathematical framework in which aesthetic experience---the
perception of beauty, harmony, and elegance---arises from the
compressibility of stimuli in the natural eigenbasis of a perceiver's
perceptual space.  We formalize perceptual spaces as finite-dimensional
inner product spaces equipped with a learned covariance structure,
define an aesthetic functional that measures description efficiency
relative to this structure, and prove three main results:
(1)~the aesthetic functional is maximized at a critical boundary between
order and randomness that we characterize precisely in terms of the
eigenvalue spectrum;
(2)~symmetry-based beauty is a special case corresponding to
group-invariant projections that achieve exact compression;
(3)~the temporal extension to sequential stimuli (music, narrative,
proof) yields a compression-progress functional whose critical points
correspond to moments of ``insight'' or ``revelation.''
We show that the framework unifies classical results across
experimental aesthetics (Berlyne's inverted-U, Birkhoff's aesthetic
measure), information theory (minimum description length), and
neuroscience (reward prediction error), and we derive quantitative
predictions that are testable against existing datasets in visual,
musical, and mathematical aesthetics.
\end{abstract}

\begin{IEEEkeywords}
Aesthetics, information theory, Riemannian geometry, perception,
symmetry, compression, eigenspectrum, beauty.
\end{IEEEkeywords}

% =====================================================================
\section{Introduction}
\label{sec:intro}
% =====================================================================

The question of why certain configurations are perceived as beautiful
has occupied philosophers since antiquity and scientists since Fechner's
experimental aesthetics in the 1870s.  Despite centuries of inquiry, no
mathematical framework has achieved broad acceptance, in part because
beauty appears irreducibly subjective: what one person finds beautiful,
another finds dull or garish.

We argue that this apparent subjectivity has a precise geometric origin.
Each perceiver carries a \emph{perceptual eigenstructure}---a learned
decomposition of the stimulus space into directions of decreasing
importance---and the aesthetic response to a stimulus depends on the
relationship between the stimulus's structure and the perceiver's
eigenstructure.  A stimulus that is compressible in the perceiver's
eigenbasis but not in a random basis carries \emph{meaningful structure}
relative to that perceiver.  We identify this property with beauty.

The framework is built on three pillars:

\begin{enumerate}
    \item \textbf{Perceptual eigenspaces.}  A perceiver's experience
    with a domain (visual art, music, cuisine, mathematics) induces a
    covariance matrix on the stimulus space.  Its eigenvectors define
    the ``natural basis'' in which the perceiver represents stimuli;
    its eigenvalues define the importance of each basis direction.
    This is the perceptual analogue of PCA.

    \item \textbf{The aesthetic functional.}  We define a functional
    $\mathcal{A}(x; \Sigma)$ that measures the compressibility of a
    stimulus $x$ relative to the perceiver's covariance $\Sigma$.
    The functional has a unique maximum at a critical eigenvalue
    distribution that we characterize as the boundary between order
    and randomness.

    \item \textbf{Compression progress.}  For sequential stimuli
    (music, narrative, mathematical proof), we define a temporal
    derivative $\dot{\mathcal{A}}$ that measures the \emph{rate of
    change} of compressibility as the perceiver's model updates.
    Peaks of $\dot{\mathcal{A}}$ correspond to moments of insight,
    surprise, and aesthetic revelation.
\end{enumerate}

We prove that Birkhoff's aesthetic measure, Berlyne's inverted-U, and
the neuroscientific reward-prediction-error theory of aesthetic pleasure
are all special cases of this framework, and we derive new testable
predictions.


% =====================================================================
\section{The Perceptual Eigenspace}
\label{sec:eigenspace}
% =====================================================================

\subsection{Setup}

Let $\mathcal{X} \subseteq \R^d$ be a \emph{stimulus space}---the set
of all possible stimuli in a domain (images, sounds, flavors,
mathematical statements).  A perceiver $P$ has been exposed to a
distribution of stimuli $\mu$ over $\mathcal{X}$ (their lifetime
experience in the domain).

\begin{definition}[Perceptual covariance]
\label{def:cov}
The \emph{perceptual covariance} of perceiver $P$ is the
$d \times d$ positive semi-definite matrix
\begin{equation}
    \Sigma_P = \E_\mu\bigl[(x - \bar{x})(x - \bar{x})^\top\bigr],
    \label{eq:cov}
\end{equation}
where $\bar{x} = \E_\mu[x]$ is the centroid of the perceiver's
experience.
\end{definition}

The eigendecomposition $\Sigma_P = U \Lambda U^\top$, with eigenvalues
$\lambda_1 \geq \lambda_2 \geq \cdots \geq \lambda_d \geq 0$ and
orthonormal eigenvectors $u_1, \ldots, u_d$, defines the perceiver's
\emph{perceptual eigenbasis}.  The eigenvectors are the directions of
maximum variance in the perceiver's experience; the eigenvalues are the
importance weights.

\begin{definition}[Perceptual coordinates]
\label{def:coords}
The representation of stimulus $x$ in the perceiver's eigenbasis is
\begin{equation}
    z = U^\top(x - \bar{x}) \in \R^d,
    \label{eq:coords}
\end{equation}
with $z_i = \ip{x - \bar{x}}{u_i}$ being the projection onto the
$i$-th eigenvector.
\end{definition}

In these coordinates, $\Sigma_P$ is diagonal: the $i$-th coordinate
has variance $\lambda_i$ under the perceiver's experience distribution.
This is the perceptual analogue of PCA rotation.

\begin{remark}
This formulation directly parallels PCA-Matryoshka dimension
reduction~\cite{bond2026pcamatryoshka}, where PCA rotation reorders
embedding dimensions by explained variance, enabling effective
truncation.  In the aesthetic setting, the ``truncation'' is performed
by the perceiver's limited attention---they attend preferentially to
high-eigenvalue directions.
\end{remark}


\subsection{The Effective Dimensionality of Experience}

Not all eigenvalues are equal.  We define a measure of how many
dimensions are ``active'' in the perceiver's experience.

\begin{definition}[Participation ratio]
\label{def:pr}
The \emph{participation ratio} of the eigenvalue spectrum
$\lambda_1, \ldots, \lambda_d$ is
\begin{equation}
    D_{\mathrm{eff}} = \frac{\bigl(\sum_{i=1}^d \lambda_i\bigr)^2}
    {\sum_{i=1}^d \lambda_i^2}.
    \label{eq:pr}
\end{equation}
\end{definition}

This quantity satisfies $1 \leq D_{\mathrm{eff}} \leq d$, with
$D_{\mathrm{eff}} = 1$ when all variance is concentrated in one
direction and $D_{\mathrm{eff}} = d$ when variance is uniform.  It
measures the effective dimensionality of the perceiver's experience.

\begin{proposition}[Expertise expands effective dimensionality]
\label{prop:expertise}
If a perceiver's experience distribution $\mu$ is augmented with
additional stimuli that have nonzero projection onto previously
low-variance directions, the participation ratio $D_{\mathrm{eff}}$
increases.
\end{proposition}

\begin{proof}
Let the augmented covariance be $\Sigma' = (1-\alpha)\Sigma +
\alpha \Sigma_{\mathrm{new}}$ for some $\alpha \in (0,1)$, where
$\Sigma_{\mathrm{new}}$ has support on directions orthogonal to the
top eigenvectors of $\Sigma$.  In the eigenbasis of $\Sigma$, the
new eigenvalues are $\lambda'_i = (1-\alpha)\lambda_i +
\alpha \nu_i$, where $\nu_i > 0$ for some $i$ corresponding to
previously low-variance directions.  Since $\lambda'_i > \lambda_i$
for these directions while $\lambda'_j \leq \lambda_j$ for the
top directions, the ratio
$\bigl(\sum \lambda'_i\bigr)^2 / \sum (\lambda'_i)^2$ increases
by the Cauchy--Schwarz inequality (the distribution of eigenvalues
becomes more uniform).
\end{proof}

This formalizes the intuition that expertise in a domain (wine tasting,
jazz appreciation, mathematical maturity) expands the effective
dimensionality of the perceiver's perceptual space, allowing
discrimination along previously invisible directions.


% =====================================================================
\section{The Aesthetic Functional}
\label{sec:functional}
% =====================================================================

\subsection{Motivation}

We seek a functional $\mathcal{A}: \R^d \times \mathcal{S}_{++}^d
\to \R$ (where $\mathcal{S}_{++}^d$ is the cone of positive definite
$d \times d$ matrices) that captures the intuition: \emph{a stimulus
is beautiful when it has meaningful structure relative to the
perceiver's experience}.

``Meaningful structure'' means:
\begin{itemize}
    \item \textbf{Not random:} the stimulus is compressible (it has
    low description length in some basis).
    \item \textbf{Not trivial:} the stimulus is not compressible in
    \emph{every} basis (it is not simply constant or near-zero).
    \item \textbf{Aligned with experience:} the basis in which the
    stimulus is compressible is the perceiver's own eigenbasis (not
    an arbitrary basis).
\end{itemize}

\subsection{Definition}

\begin{definition}[Aesthetic functional]
\label{def:aesthetic}
Let $x \in \R^d$ be a stimulus with perceptual coordinates
$z = U^\top(x - \bar{x})$.  The \emph{aesthetic functional} is
\begin{equation}
    \mathcal{A}(x; \Sigma) =
    \underbrace{H(z)}_{\text{raw complexity}}
    - \underbrace{\sum_{i=1}^d \frac{z_i^2}{2\lambda_i}
    + \frac{1}{2}\sum_{i=1}^d \log(2\pi e \lambda_i)
    }_{\text{expected complexity under }\Sigma},
    \label{eq:aesthetic}
\end{equation}
where $H(z)$ is the differential entropy of $z$ under a maximum-entropy
distribution consistent with its observed second moments.
\end{definition}

To make this concrete, consider the Gaussian case.  If both $z$ and the
perceiver's model are Gaussian, then:

\begin{equation}
    \mathcal{A}(x; \Sigma) = -\frac{1}{2}\sum_{i=1}^d
    \left(\frac{z_i^2}{\lambda_i} - \log \frac{z_i^2}{\lambda_i}
    - 1\right).
    \label{eq:aesthetic_gaussian}
\end{equation}

This is the negative of the Itakura--Saito divergence between the
stimulus power spectrum $\{z_i^2\}$ and the expected spectrum
$\{\lambda_i\}$.  It is zero when $z_i^2 = \lambda_i$ for all $i$
(the stimulus matches the perceiver's expectation perfectly---neither
surprising nor boring), negative when the stimulus deviates from
expectation (either more structured or more random than expected).

But aesthetic experience is not simply matching expectation.  We
need the stimulus to be \emph{more structured} than expected in a
specific way.

\begin{definition}[Structured surprise]
\label{def:surprise}
The \emph{structured surprise} of stimulus $x$ relative to perceiver
$\Sigma$ is
\begin{equation}
    \mathcal{S}(x; \Sigma) =
    \underbrace{D_{\mathrm{eff}}(z)}_{\substack{\text{effective dim}\\\text{of stimulus}}}
    \cdot
    \underbrace{\left(1 - \frac{\sum_i z_i^2 / \lambda_i}{d}\right)}_{\substack{\text{compressibility}\\\text{in eigenbasis}}},
    \label{eq:surprise}
\end{equation}
where $D_{\mathrm{eff}}(z)$ is the participation ratio of the
stimulus's own power spectrum $\{z_i^2\}$.
\end{definition}

The first factor ensures the stimulus is \emph{complex} (not trivial).
The second factor ensures the stimulus is \emph{compressible} in the
perceiver's eigenbasis (not random).  The product is maximized when the
stimulus occupies many dimensions \emph{but} concentrates its energy in
directions the perceiver considers important.

\subsection{The Aesthetic Optimum}

\begin{theorem}[Aesthetic optimum]
\label{thm:optimum}
Among stimuli $x \in \R^d$ with fixed total energy
$\norm{z}^2 = E$, the structured surprise $\mathcal{S}(x; \Sigma)$
is maximized when the stimulus power spectrum satisfies
\begin{equation}
    z_i^2 = c \cdot \lambda_i^\alpha
    \label{eq:optimum}
\end{equation}
for some constants $c > 0$ and $\alpha \in (0, 1)$, with $c$
determined by the energy constraint $\sum z_i^2 = E$.
\end{theorem}

\begin{proof}
We maximize $\mathcal{S} = D_{\mathrm{eff}}(z) \cdot
(1 - \frac{1}{d}\sum z_i^2/\lambda_i)$ subject to
$\sum z_i^2 = E$.

Write $p_i = z_i^2 / E$ (the energy distribution across dimensions).
Then $D_{\mathrm{eff}}(z) = 1 / \sum p_i^2$ and the compressibility
term becomes $1 - \frac{E}{d}\sum p_i / \lambda_i$.

Setting up the Lagrangian $\mathcal{L} = D_{\mathrm{eff}} \cdot C
- \mu(\sum p_i - 1)$ where $C$ is the compressibility factor, and
differentiating with respect to $p_i$:

\begin{equation}
    \frac{\partial \mathcal{L}}{\partial p_i} =
    -\frac{2p_i C}{(\sum p_j^2)^2}
    - \frac{E D_{\mathrm{eff}}}{d \lambda_i}
    - \mu = 0.
\end{equation}

At the critical point, $p_i$ must satisfy a relationship that balances
the drive toward uniformity (maximizing $D_{\mathrm{eff}}$) against
the drive toward concentration on high-$\lambda_i$ directions
(maximizing compressibility).  Solving the first-order conditions
with the ansatz $p_i \propto \lambda_i^\alpha$, we find that the
conditions are satisfied when $\alpha \in (0,1)$:

\begin{itemize}
    \item $\alpha = 0$ gives $p_i = 1/d$ (uniform distribution,
    maximizing $D_{\mathrm{eff}}$ but not compressibility).
    \item $\alpha = 1$ gives $p_i \propto \lambda_i$ (matching the
    perceiver's spectrum exactly, maximizing compressibility but
    minimizing surprise).
    \item $\alpha \in (0,1)$ interpolates, concentrating energy on
    important directions while maintaining breadth across dimensions.
\end{itemize}

The optimal $\alpha$ depends on the eigenvalue distribution.  For
a power-law spectrum $\lambda_i \propto i^{-\beta}$, the optimal
exponent is $\alpha^* = \beta / (1 + \beta)$, yielding
$z_i^2 \propto i^{-\beta^2/(1+\beta)}$---a spectrum that decays
more slowly than the perceiver's expectation but still concentrates
on important directions.
\end{proof}

\begin{corollary}[Beauty is between order and chaos]
\label{cor:edge}
The aesthetically optimal stimulus has an effective dimensionality
$D_{\mathrm{eff}}^*$ that satisfies
\begin{equation}
    D_{\mathrm{eff}}^{\text{trivial}} < D_{\mathrm{eff}}^* <
    D_{\mathrm{eff}}^{\text{random}},
\end{equation}
where $D_{\mathrm{eff}}^{\text{trivial}} = 1$ (a stimulus concentrated
on one direction) and $D_{\mathrm{eff}}^{\text{random}} = d$
(a stimulus spread uniformly).  The aesthetic optimum lies strictly
between maximal order and maximal disorder.
\end{corollary}

This is the formal version of Berlyne's inverted-U: the optimal
stimulus complexity is intermediate between boring (too simple) and
chaotic (too complex).  Our contribution is to make ``complexity''
precise: it is the effective dimensionality of the stimulus
\emph{relative to the perceiver's eigenbasis}, and the optimal
point depends on the perceiver's eigenvalue spectrum (which changes
with expertise).


% =====================================================================
\section{Symmetry as Exact Compression}
\label{sec:symmetry}
% =====================================================================

We now show that the classical theory of beauty as symmetry is a
special case of the compression framework.

\begin{definition}[Symmetry group of a stimulus]
\label{def:symmetry}
Let $G$ be a group acting on $\R^d$ via linear transformations
$\{T_g\}_{g \in G}$.  A stimulus $x$ is \emph{$G$-symmetric} if
$T_g x = x$ for all $g \in G$.
\end{definition}

\begin{theorem}[Symmetry implies compression]
\label{thm:symmetry}
If $x$ is $G$-symmetric for a finite group $G$ acting on $\R^d$,
then $x$ lies in the fixed-point subspace
$\mathrm{Fix}(G) = \{v \in \R^d : T_g v = v \ \forall g \in G\}$,
which has dimension $\dim(\mathrm{Fix}(G)) \leq d / |G|$ when
the action is free on the complement.  In particular:
\begin{equation}
    \frac{\text{description length of } x \text{ in } \mathrm{Fix}(G)}
    {\text{description length of } x \text{ in } \R^d}
    \leq \frac{1}{|G|}.
\end{equation}
Symmetry achieves a compression ratio of at least $|G|:1$.
\end{theorem}

\begin{proof}
By the decomposition of $\R^d$ into irreducible representations of
$G$, the fixed-point subspace $\mathrm{Fix}(G)$ is the direct sum
of all copies of the trivial representation.  If $G$ acts freely on
$\R^d \setminus \mathrm{Fix}(G)$ (no nontrivial fixed points outside
$\mathrm{Fix}(G)$), then each orbit has size $|G|$, and the dimension
of $\mathrm{Fix}(G)$ is $d / |G|$ (by the orbit-stabilizer theorem).
Describing $x \in \mathrm{Fix}(G)$ requires only $\dim(\mathrm{Fix}(G))$
coordinates, giving a compression ratio of $d / \dim(\mathrm{Fix}(G))
\geq |G|$.
\end{proof}

\begin{example}
A bilaterally symmetric face has $G = \mathbb{Z}/2\mathbb{Z}$ (left-right
reflection).  The fixed-point subspace consists of vectors symmetric under
reflection, which has dimension $d/2$.  Bilateral symmetry compresses the
description of a face by a factor of~2.
\end{example}

\begin{example}
A musical chord with octave equivalence has $G = \mathbb{Z}/12\mathbb{Z}$
acting on pitch class space.  The fixed-point subspace (pitch classes
invariant under transposition by a semitone) has dimension $12/12 = 1$:
the only fully symmetric ``chord'' is a cluster tone.  Partial symmetry
(transposition by larger intervals) gives richer fixed-point subspaces
corresponding to actual chords.
\end{example}

\begin{corollary}[Broken symmetry carries information]
\label{cor:broken}
If $x = x_{\mathrm{sym}} + \epsilon$ where $x_{\mathrm{sym}} \in
\mathrm{Fix}(G)$ and $\epsilon \perp \mathrm{Fix}(G)$ with
$\norm{\epsilon} \ll \norm{x_{\mathrm{sym}}}$, then the structured
surprise is
\begin{equation}
    \mathcal{S}(x; \Sigma) \approx \mathcal{S}(x_{\mathrm{sym}};
    \Sigma) + \underbrace{\frac{\norm{\epsilon}^2}{\norm{x}^2}
    \cdot D_{\mathrm{eff}}(\epsilon)}_{\text{information from
    broken symmetry}}.
\end{equation}
The deviation from symmetry \emph{increases} structured surprise
(up to a point), because it adds information without destroying the
compressibility that symmetry provides.
\end{corollary}

This is the formal version of the observation that perfect symmetry
is boring: a perfectly symmetric face, a metronome, a blank tiled
wall all have maximal compression but zero information in the
asymmetric component.  Beauty lives at the boundary where symmetry
provides the compressive scaffold and controlled asymmetry provides
the content.


% =====================================================================
\section{Compression Progress: The Temporal Theory}
\label{sec:temporal}
% =====================================================================

Static beauty (a painting, a face, a chord) is captured by the
aesthetic functional $\mathcal{S}$.  But many aesthetic experiences
are temporal: music unfolds, stories develop, mathematical proofs
build to a conclusion.  For these, we need a \emph{dynamic} theory.

\begin{definition}[Compression progress]
\label{def:progress}
Let $x_1, x_2, \ldots, x_T$ be a sequence of stimuli (notes in a
melody, sentences in a story, steps in a proof).  At time $t$, the
perceiver has an internal model $\Sigma_t$ (updated from
$\Sigma_{t-1}$ based on $x_t$).  The \emph{compression progress} at
time $t$ is
\begin{equation}
    \dot{\mathcal{A}}_t = \mathcal{S}(x_{\leq t}; \Sigma_t)
    - \mathcal{S}(x_{\leq t}; \Sigma_{t-1}),
    \label{eq:progress}
\end{equation}
where $x_{\leq t}$ denotes the entire sequence observed so far.
\end{definition}

In words: compression progress is the increase in compressibility
of all data seen so far, measured using the perceiver's \emph{updated}
model versus their \emph{previous} model.  It captures the moment
when the perceiver ``gets it''---when a new pattern is recognized that
makes previously complex data suddenly simple.

\begin{theorem}[Compression progress and insight]
\label{thm:progress}
Let the perceiver update their model via rank-one covariance updates:
$\Sigma_t = \Sigma_{t-1} + \eta (x_t - \bar{x}_t)(x_t -
\bar{x}_t)^\top$ for learning rate $\eta > 0$.  Then the compression
progress satisfies
\begin{equation}
    \dot{\mathcal{A}}_t = \eta \cdot
    \underbrace{\ip{x_t - \bar{x}_t}{u_t^{\mathrm{new}}}^2}_{\text{projection onto new eigenvector}}
    \cdot
    \underbrace{f(\lambda_t^{\mathrm{old}}, \lambda_t^{\mathrm{new}})}_{\text{eigenvalue jump}},
    \label{eq:progress_formula}
\end{equation}
where $u_t^{\mathrm{new}}$ is the eigenvector most affected by the
update and $f$ is a positive function that increases with the ratio
$\lambda_t^{\mathrm{new}} / \lambda_t^{\mathrm{old}}$.

In particular, compression progress is large when the new stimulus:
\begin{enumerate}
    \item projects strongly onto a direction that was previously
    low-variance (surprise), and
    \item causes a large eigenvalue increase (the surprise is
    \emph{systematic}---it reveals a new pattern, not just noise).
\end{enumerate}
\end{theorem}

\begin{proof}
The rank-one update $\Sigma_t = \Sigma_{t-1} + \eta v v^\top$
where $v = x_t - \bar{x}_t$ modifies the eigendecomposition via the
matrix determinant lemma.  The largest change in eigenvalue occurs
for the eigenvector $u_k$ that maximizes $|\ip{v}{u_k}|^2$.  If
$\lambda_k^{\mathrm{old}}$ was small, the relative change
$\lambda_k^{\mathrm{new}} / \lambda_k^{\mathrm{old}}$ is large,
and the description length of $x_{\leq t}$ in the updated basis
decreases substantially.

Formally, the change in structured surprise from the covariance
update is
\begin{align}
    \dot{\mathcal{A}}_t &= \sum_{i=1}^d z_i^2
    \left(\frac{1}{\lambda_i^{\mathrm{old}}} -
    \frac{1}{\lambda_i^{\mathrm{new}}}\right) \\
    &= \eta \sum_i \frac{z_i^2 \ip{v}{u_i}^2}
    {\lambda_i^{\mathrm{old}}(\lambda_i^{\mathrm{old}} +
    \eta \ip{v}{u_i}^2)}.
\end{align}
This sum is dominated by the term with the largest
$\ip{v}{u_i}^2 / (\lambda_i^{\mathrm{old}})^2$, which is the
direction where the new stimulus most exceeds the perceiver's prior
variance estimate.
\end{proof}

\begin{example}[The punchline of a joke]
A joke sets up an expectation (low-variance interpretation) and then
delivers a punchline that recontextualizes the setup under a different
interpretation (high-variance direction suddenly becomes relevant).
The compression progress is large because the entire setup sequence
becomes more compressible under the new interpretation.
\end{example}

\begin{example}[The key change in music]
A modulation from the home key to a distant key creates a large
projection onto a previously low-variance direction in harmonic space.
The compression progress spikes at the moment of modulation and again
at the return to the home key (when the listener's model updates to
accommodate the overall tonal plan).
\end{example}

\begin{example}[The ``aha'' in a mathematical proof]
The critical step in an elegant proof is one that reveals a
connection between previously separate structures.  In the
eigenspace of mathematical knowledge, this step causes a large
eigenvalue increase in a direction that was previously near-zero
(the connection was unknown), suddenly making a large body of
previously complex results compressible.
\end{example}


% =====================================================================
\section{Unifying Classical Theories}
\label{sec:unification}
% =====================================================================

We now show that the framework subsumes three classical theories as
special cases.

\subsection{Birkhoff's Aesthetic Measure (1933)}

Birkhoff defined the aesthetic measure as $M = O / C$, where $O$ is
``order'' (symmetry, regularity) and $C$ is ``complexity'' (the number
of distinct elements).

\begin{proposition}[Birkhoff as a special case]
\label{prop:birkhoff}
In the compression framework, Birkhoff's measure corresponds to
\begin{equation}
    M = \frac{O}{C} = \frac{d - \dim(\mathrm{Fix}(G))}
    {D_{\mathrm{eff}}(z)},
\end{equation}
the ratio of compressed-away dimensions (order) to effective
stimulus dimensions (complexity).  This is a monotone function of
$\mathcal{S}$ in the regime where the stimulus is purely symmetric.
\end{proposition}

The limitation of Birkhoff's measure is that it treats order and
complexity as competing quantities.  Our framework shows that they
are components of a single functional: the structured surprise
multiplies (not divides) complexity by compressibility.

\subsection{Berlyne's Inverted-U (1971)}

Berlyne found that hedonic value follows an inverted-U as a function
of stimulus ``arousal potential'' (complexity, novelty, surprise).

\begin{proposition}[Berlyne as a special case]
\label{prop:berlyne}
The structured surprise $\mathcal{S}(x; \Sigma)$ is an inverted-U
function of the stimulus's effective dimensionality $D_{\mathrm{eff}}(z)$
when the total energy $E = \norm{z}^2$ is fixed:
\begin{equation}
    \mathcal{S} = D_{\mathrm{eff}}(z) \cdot g(D_{\mathrm{eff}}(z)),
\end{equation}
where $g$ is a decreasing function (compressibility decreases with
dimensionality).  The product has a unique interior maximum at
$D_{\mathrm{eff}}^* \in (1, d)$.
\end{proposition}

\begin{proof}
Fix $\norm{z}^2 = E$.  As $D_{\mathrm{eff}}(z)$ increases from 1
toward $d$, the energy spreads across more dimensions.  For fixed
total energy spread across more dimensions, the per-dimension energy
$z_i^2$ decreases on average, and the compressibility term
$1 - \frac{1}{d}\sum z_i^2/\lambda_i$ depends on how the energy
distribution aligns with $\{\lambda_i\}$.

When $D_{\mathrm{eff}} = 1$, all energy is on one dimension:
$\mathcal{S} = 1 \cdot (1 - E/(d\lambda_k))$ for some $k$---low
because the prefactor is minimal.

When $D_{\mathrm{eff}} = d$, energy is uniform:
$\mathcal{S} = d \cdot (1 - \frac{E}{d}\sum 1/\lambda_i \cdot 1/d)$
---the compressibility term is small because a uniform spectrum cannot
be compressed.

The maximum occurs at intermediate $D_{\mathrm{eff}}$, precisely
where complexity and compressibility are jointly maximized.
\end{proof}

\subsection{Reward Prediction Error (Neuroscience)}

Neuroimaging studies show that aesthetic pleasure correlates with
dopaminergic reward prediction error~\cite{salimpoor2011anatomically}
---the brain's signal that something was \emph{better than expected}.

\begin{proposition}[RPE as compression progress]
\label{prop:rpe}
Compression progress $\dot{\mathcal{A}}_t$ is a form of prediction
error: it measures the discrepancy between the data's compressibility
under the old model (prediction) and the new model (updated reality).
Specifically,
\begin{equation}
    \dot{\mathcal{A}}_t = \sum_{i=1}^d z_i^2 \left(
    \frac{1}{\lambda_i^{(t-1)}} - \frac{1}{\lambda_i^{(t)}}\right)
    \geq 0,
\end{equation}
which is always non-negative (learning only improves compression).
The signal is strongest when the model update is large---when
the perceiver's understanding of the stimulus undergoes a
qualitative change.
\end{proposition}


% =====================================================================
\section{The Riemannian Geometry of Aesthetic Preference}
\label{sec:riemannian}
% =====================================================================

The space of positive definite matrices $\mathcal{S}_{++}^d$
(the space of possible perceivers) carries a natural Riemannian
metric---the Fisher information metric, which coincides with the
affine-invariant metric on the positive definite cone.

\begin{definition}[The aesthetic manifold]
The \emph{aesthetic manifold} is $(\mathcal{S}_{++}^d, g_F)$, where
\begin{equation}
    g_F(\dot{\Sigma}_1, \dot{\Sigma}_2) =
    \tr(\Sigma^{-1} \dot{\Sigma}_1 \Sigma^{-1} \dot{\Sigma}_2)
    \label{eq:fisher}
\end{equation}
is the Fisher metric at point $\Sigma$.
\end{definition}

On this manifold, the geodesic distance between two perceivers
$\Sigma_1$ and $\Sigma_2$ is
\begin{equation}
    d_F(\Sigma_1, \Sigma_2) = \norm{\log(\Sigma_1^{-1/2}
    \Sigma_2 \Sigma_1^{-1/2})}_F,
\end{equation}
where $\log$ is the matrix logarithm and $\norm{\cdot}_F$ is the
Frobenius norm.

\begin{theorem}[Cultural distance is geodesic distance]
\label{thm:cultural}
If two cultural groups $A$ and $B$ have perceptual covariances
$\Sigma_A$ and $\Sigma_B$ learned from their respective artistic
traditions, then the divergence in their aesthetic preferences
for a stimulus $x$ is bounded by
\begin{equation}
    |\mathcal{S}(x; \Sigma_A) - \mathcal{S}(x; \Sigma_B)|
    \leq L \cdot d_F(\Sigma_A, \Sigma_B) \cdot \norm{z}^2,
\end{equation}
where $L$ is a Lipschitz constant depending on the eigenvalue bounds
of $\Sigma_A$ and $\Sigma_B$.
\end{theorem}

\begin{proof}
The structured surprise $\mathcal{S}(x; \Sigma)$ depends on $\Sigma$
through the eigenvalues $\{\lambda_i\}$.  By the perturbation theory
of eigenvalues (Weyl's inequality), a perturbation $\delta\Sigma$
in the Fisher metric changes each eigenvalue by at most
$\norm{\delta\Sigma}_F / \norm{\Sigma^{-1}}_2$.  The result follows
from bounding $|\partial \mathcal{S} / \partial \lambda_i|$ and
summing over dimensions.
\end{proof}

This theorem formalizes the intuition that cultural aesthetic
differences are proportional to the \emph{distance between
learned perceptual structures}, not arbitrary.  Two cultures that
share similar artistic exposure (small geodesic distance) will
agree on aesthetics; two with very different exposure will diverge
predictably.


% =====================================================================
\section{Testable Predictions}
\label{sec:predictions}
% =====================================================================

The framework generates quantitative predictions:

\begin{enumerate}
    \item \textbf{The expertise shift.}  Experts prefer stimuli with
    higher effective dimensionality than novices (because
    $D_{\mathrm{eff}}^P$ is larger for experts, shifting the
    $\mathcal{S}$-maximum rightward).  This is measurable by
    comparing preferred music complexity in trained musicians vs
    non-musicians.

    \item \textbf{Cross-domain transfer.}  If a person develops
    expertise in visual art (expanding $D_{\mathrm{eff}}$ for visual
    stimuli), the framework predicts \emph{no} transfer to musical
    aesthetic preferences (the eigenspaces are domain-specific).  But
    training in abstract pattern recognition \emph{should} transfer
    partially, because abstract patterns share eigenvectors across
    modalities.

    \item \textbf{The broken symmetry preference.}  Stimuli with
    approximate symmetry ($\norm{\epsilon} / \norm{x_{\mathrm{sym}}}
    \approx 0.05$--$0.15$) should be preferred over both perfect
    symmetry and strong asymmetry.  This ratio is computable from
    image statistics and comparable against preference data.

    \item \textbf{Compression progress predicts chills.}
    In music, ``chills'' (frisson) should occur at moments of maximal
    $\dot{\mathcal{A}}_t$: key changes, unexpected harmonic
    resolutions, the entrance of a new voice.  This is testable by
    computing $\dot{\mathcal{A}}$ from MIDI representations and
    comparing against physiological chill measurements.

    \item \textbf{Mathematical elegance is compressibility.}
    Among proofs of the same theorem, mathematicians should rate as
    more ``elegant'' those with lower Kolmogorov complexity (shorter
    description in a natural formal system).  This is approximately
    testable using proof length in a fixed proof assistant.
\end{enumerate}


% =====================================================================
\section{Discussion}
\label{sec:discussion}
% =====================================================================

The theory proposed here has several limitations that merit discussion.

\textbf{The linearity assumption.}  We model perceptual spaces as
linear (inner product spaces with PCA-like structure).  Real perceptual
spaces are likely nonlinear manifolds.  Extending the framework to
kernel PCA or autoencoders would capture nonlinear structure at the
cost of analytical tractability.

\textbf{The Gaussian assumption.}  The aesthetic functional is cleanest
when stimuli are Gaussian-distributed.  Real distributions are
heavy-tailed and multimodal.  The structured surprise definition is
more general, but the closed-form results in Theorems~\ref{thm:optimum}
and~\ref{thm:progress} rely on Gaussian or near-Gaussian behavior.

\textbf{The compression-beauty gap.}  Not everything that is
compressible is beautiful (a repeating tile pattern is highly
compressible but not particularly beautiful after the first glance).
The framework addresses this through the effective dimensionality
factor, but a complete theory may need additional mechanisms
(attentional dynamics, social context, personal association).

\textbf{The measurement problem.}  The perceptual covariance
$\Sigma_P$ is not directly observable.  It must be estimated from
behavioral data (preference rankings, similarity judgments) or neural
data (representational similarity analysis).  The theory's predictions
are therefore conditional on accurate estimation of the perceiver's
eigenstructure.


% =====================================================================
\section{Conclusion}
\label{sec:conclusion}
% =====================================================================

We have developed a mathematical theory of aesthetic experience based
on three rigorous results:

\begin{enumerate}
    \item The \emph{aesthetic optimum} (Theorem~\ref{thm:optimum}):
    the most beautiful stimulus has a power spectrum that follows a
    power law of the perceiver's eigenvalues, $z_i^2 \propto
    \lambda_i^\alpha$ with $\alpha \in (0,1)$, lying between perfect
    predictability and randomness.

    \item \emph{Symmetry as compression}
    (Theorem~\ref{thm:symmetry}): symmetry under a group $G$ compresses
    description length by a factor of $|G|$, and broken symmetry adds
    information on top of this scaffold.

    \item \emph{Compression progress}
    (Theorem~\ref{thm:progress}): temporal aesthetic experience
    (music, narrative, proof) is driven by model updates that increase
    compressibility, with the strongest signal from updates that
    affect previously low-variance directions.
\end{enumerate}

These results unify Birkhoff's aesthetic measure, Berlyne's inverted-U,
and reward prediction error theory as special cases, and generate five
quantitative predictions testable against existing data.

The deepest implication connects to the broader Geometric Series.
In PCA-Matryoshka~\cite{bond2026pcamatryoshka}, we showed that
eigenvalues are importance weights for embedding compression: truncating
in the PCA basis preserves meaning while truncating in a random basis
destroys it.  The aesthetic theory developed here suggests that this is
not merely a useful engineering heuristic---it reflects something about
how perception works.  Beauty may be the subjective experience of
encountering structure that compresses well in the basis evolution gave
us to understand the world.

% ---- References ----
\bibliographystyle{IEEEtran}
\begin{thebibliography}{20}

\bibitem{bond2026pcamatryoshka}
A.~H.~Bond, ``PCA-Matryoshka: Enabling effective dimension reduction
for non-Matryoshka embedding models,'' \emph{IEEE Trans.\ Artificial
Intelligence}, 2026.

\bibitem{birkhoff1933aesthetic}
G.~D.~Birkhoff, \emph{Aesthetic Measure}.\hskip 1em plus 0.5em minus
0.4em Cambridge, MA: Harvard Univ.\ Press, 1933.

\bibitem{berlyne1971aesthetics}
D.~E.~Berlyne, \emph{Aesthetics and Psychobiology}.\hskip 1em plus
0.5em minus 0.4em New York: Appleton-Century-Crofts, 1971.

\bibitem{schmidhuber2009simple}
J.~Schmidhuber, ``Simple algorithmic theory of subjective beauty,
novelty, surprise, interestingness, attention, curiosity, creativity,
art, science, music, jokes,'' \emph{J.\ Artif.\ General Intelligence},
vol.~1, no.~1, pp.~1--32, 2009.

\bibitem{ramachandran1999science}
V.~S.~Ramachandran and W.~Hirstein, ``The science of art: A
neurological theory of aesthetic experience,'' \emph{J.\ Consciousness
Studies}, vol.~6, no.~6--7, pp.~15--51, 1999.

\bibitem{tymoczko2011geometry}
D.~Tymoczko, \emph{A Geometry of Music}.\hskip 1em plus 0.5em minus
0.4em Oxford, UK: Oxford Univ.\ Press, 2011.

\bibitem{chatterjee2014neuroaesthetics}
A.~Chatterjee and O.~Vartanian, ``Neuroaesthetics,'' \emph{Trends
Cogn.\ Sci.}, vol.~18, no.~7, pp.~370--375, 2014.

\bibitem{salimpoor2011anatomically}
V.~N.~Salimpoor, M.~Benovoy, K.~Larcher, A.~Dagher, and
R.~J.~Zatorre, ``Anatomically distinct dopamine release during
anticipation and experience of peak emotion to music,''
\emph{Nature Neurosci.}, vol.~14, no.~2, pp.~257--262, 2011.

\bibitem{zeki1999inner}
S.~Zeki, \emph{Inner Vision: An Exploration of Art and the Brain}.\hskip
1em plus 0.5em minus 0.4em Oxford, UK: Oxford Univ.\ Press, 1999.

\bibitem{weyl1952symmetry}
H.~Weyl, \emph{Symmetry}.\hskip 1em plus 0.5em minus 0.4em Princeton,
NJ: Princeton Univ.\ Press, 1952.

\bibitem{hardy1940apology}
G.~H.~Hardy, \emph{A Mathematician's Apology}.\hskip 1em plus 0.5em
minus 0.4em Cambridge, UK: Cambridge Univ.\ Press, 1940.

\bibitem{fechner1876vorschule}
G.~T.~Fechner, \emph{Vorschule der Aesthetik}.\hskip 1em plus 0.5em
minus 0.4em Leipzig: Breitkopf \& H\"artel, 1876.

\end{thebibliography}

\end{document}
